\subsection{Basis-function ablation}
\label{sec:basis-results}

We now fix the Tsit5 solver and compare the seven basis families defined in
Section~\ref{sec:bases}, starting with their conditioning and locality.

Figure~\ref{fig:basis-catalog} draws the seven families from the code
itself and reports the 2-norm condition number $\kappa_2(\Psi)$ of the
$400\times G$ design matrix $\Psi_{ng}=B_g(u_n)$ on a uniform sample of
$[-1,1]$. This is the quantity that governs how well-posed fitting the
coefficients $c_g$ is. At $G=5$ Chebyshev ($2.5$), Lagrange ($3.7$) and
B-spline ($5.5$) are well conditioned, the three kernels sit at $8$--$34$,
and Newton is worst at $54$. The growth with $G$ separates them further: at
$G=12$ Chebyshev and B-spline are still below $6$, RBF and Lagrange reach
$69$ and $63$, and Newton reaches $2.4\times10^4$, growing exponentially,
because the nested products $\prod_{j<g}(u-z_j)$ become nearly collinear
columns.

\begin{figure*}[t]
\centering
\includegraphics[width=\textwidth]{figures/analysis/basis_catalog.pdf}
\caption{The seven implemented basis families on the $G=5$ grid (vertical
lines are the knots $z_g$), with the $g=4$ function highlighted and the
2-norm condition number of the design matrix in each title. The blue band is
the interval $u\in[0.25,1]$ that first-layer inputs actually reach on the
Lotka-Volterra orbit. The last panel shows $\kappa_2$ against grid size;
dashed and dotted grey lines are IQF and RSWAF.}
\label{fig:basis-catalog}
\end{figure*}

\paragraph{How much of the grid the data uses.} Lotka-Volterra's states are
positive, between $0.26$ and $6.9$ on the orbit, so after the $\tanh$
normaliser the first layer sees only $u\in[0.25,1.0]$, the right
$37.5\%$ of its grid. Measured on the trained RBF model, the two leftmost
basis functions ($z_1=-1$, $z_2=-0.5$) carry $0.4\%$ and $3.6\%$ of the
first layer's total basis activation; for the B-spline the leftmost carries
exactly zero and the next $1.7\%$. The second layer, whose inputs are the
signed hidden activations, uses its grid evenly (every basis function
$18$--$26\%$). Roughly a third of the first layer's 100 spline coefficients
are therefore nearly inert on this system. Figure~\ref{fig:locality} shows
what locality means against that data interval: raising the coefficient of
$z_2$ changes the spline edge by at most $0.021\delta$ on the data, the RBF
edge by $0.32\delta$, and the Chebyshev edge by the full $\delta$.

\begin{figure*}[t]
\centering
\includegraphics[width=0.85\textwidth]{figures/tikz/local_vs_global_basis.pdf}
\caption{Locality of one coefficient update, $c_2\leftarrow c_2+\delta$,
for three basis families (uniform-knot cubic spline shown for clarity). Since
$\phi$ is linear in $\mathbf{c}$, the change is exactly $\delta\,\psi_2$
(shaded). Compact support confines it to its knot span; the RBF's Gaussian
tail still reaches the data interval (blue); a global polynomial changes the
edge everywhere.}
\label{fig:locality}
\end{figure*}

\begin{table*}[t]
\centering
\footnotesize
\renewcommand{\arraystretch}{1.1}
\caption{Basis ablation (fixed Tsit5 solver, $G=5$ grid points, 240
parameters, $10^4$ epochs). $\kappa_2$ is the design-matrix condition number
of Figure~\ref{fig:basis-catalog}.}
\label{tab:basis-ablation}
\begin{tabularx}{\textwidth}{@{}Y C{2.3cm} C{2.3cm} C{1.7cm} C{1.6cm} C{1.2cm} C{1.2cm} C{1.6cm}@{}}
\toprule
\hdrrow \thh{Basis function} & \thh{Training MSE} & \thh{Extrap.\ MSE} & \thh{$R^2$} & \thh{Rel.\ $L_2$} & \thh{$\hat L$} & \thh{$\kappa_2$} & \thh{Wall-clock}\\
\midrule
Cubic B-spline & $\mathbf{7.20\times10^{-5}}$ & $\mathbf{8.38\times10^{-5}}$ & $\mathbf{0.99997}$ & $\mathbf{0.32\%}$ & $8.44$ & $5.5$ & 6680 s\\
Gaussian RBF & $8.83\times10^{-5}$ & $8.92\times10^{-5}$ & $0.99997$ & $0.33\%$ & $6.91$ & $34.0$ & \textbf{1712 s}\\
Lagrange & $1.07\times10^{-4}$ & $3.93\times10^{-4}$ & $0.99987$ & $0.69\%$ & $6.70$ & $3.7$ & 6829 s\\
Chebyshev & $9.80\times10^{-5}$ & $4.37\times10^{-4}$ & $0.99985$ & $0.73\%$ & $6.48$ & $2.5$ & 2317 s\\
IQF & $2.08\times10^{-4}$ & $9.61\times10^{-4}$ & $0.99967$ & $1.08\%$ & $7.16$ & $10.1$ & 1824 s\\
RSWAF & $1.00\times10^{-4}$ & $1.46\times10^{-3}$ & $0.99950$ & $1.33\%$ & $8.72$ & $7.9$ & 1715 s\\
Newton & $1.47\times10^{-3}$ & $3.65\times10^{-2}$ & $0.98746$ & $6.66\%$ & $5.65$ & $54.2$ & 2187 s\\
\bottomrule
\end{tabularx}
\end{table*}

B-spline edges out RBF on accuracy after a boundary-knot fix, at $3.9\times$
RBF's wall-clock time and a visibly higher Lipschitz estimate ($8.44$ vs.\
$6.91$). Its compact local support lets it fit sharper local features than
RBF's smooth, slowly decaying support without being penalized elsewhere in
the domain. The two bases are statistically tied on error (6\% apart, inside
single-seed noise), which makes RBF the better accuracy-per-second choice.
Polynomial global-support bases (Chebyshev, Lagrange) fit the training
window well and extrapolate $4$--$5\times$ worse. This is the expected
signature of a basis whose every coefficient influences the whole domain
rather than a local region of it. Conditioning alone does not predict the
ranking: Chebyshev is the best-conditioned basis and extrapolates $5\times$
worse than RBF, the worst-conditioned of the working kernels. Locality
matters more than conditioning once $\kappa_2$ is moderate. The two
heavy-tailed kernels, IQF (algebraic $O(r^{-2})$ tails) and RSWAF, train as
well as the polynomials but extrapolate worst among the working bases
($11$--$16\times$ RBF). Newton is the outlier in both directions: it is the
worst-conditioned basis and the only one that never reaches $10^{-3}$
training loss, and its lowest-in-table Lipschitz estimate ($5.65$) reflects
underfitting, not desirable flatness. Figure~\ref{fig:basis-phase} shows the
same ranking in the phase plane.

\begin{figure*}[t]
\centering
\begin{subfigure}[b]{0.24\textwidth}
  \centering
  \includegraphics[width=\linewidth]{figures/ablation/basis_bspline_phase.png}
  \caption{Cubic B-spline}
\end{subfigure}\hfill
\begin{subfigure}[b]{0.24\textwidth}
  \centering
  \includegraphics[width=\linewidth]{figures/ablation/basis_rbf_phase.png}
  \caption{Gaussian RBF}
\end{subfigure}\hfill
\begin{subfigure}[b]{0.24\textwidth}
  \centering
  \includegraphics[width=\linewidth]{figures/ablation/basis_chebyshev_phase.png}
  \caption{Chebyshev}
\end{subfigure}\hfill
\begin{subfigure}[b]{0.24\textwidth}
  \centering
  \includegraphics[width=\linewidth]{figures/ablation/basis_newton_phase.png}
  \caption{Newton}
\end{subfigure}
\caption{Phase portraits on $[0,14]$ for four of the seven bases (Tsit5).
B-spline, RBF and Chebyshev all trace the true orbit closely at this scale;
Chebyshev's $5\times$ larger extrapolation error is phase error rather than
a visible departure. The Newton-basis model visibly misses the orbit near
its upper and right-hand extremes.}
\label{fig:basis-phase}
\end{figure*}
